\section{Literature Review}

Quantile regression describes how covariates are related to different parts of a conditional outcome distribution rather than only its mean \parencite{koenker_bassett_1978}. This distinction matters when treatment effects are heterogeneous: an intervention may shift lower and upper quantiles by different amounts even when its average effect is modest. A structural quantile treatment effect compares the same quantile of potential-outcome distributions under alternative treatment states. It is generally a distributional effect, not an individual-level effect, because the observation occupying a given rank can differ across counterfactual states. Endogeneity creates a separate problem. When realized treatment depends on latent determinants of the outcome, the conditional distribution among observations selecting a treatment value need not equal the potential-outcome distribution obtained by fixing treatment at that value. Ordinary quantile regression of the outcome on treatment and controls then estimates an observed conditional association and does not generally recover a structural quantile effect. The IVQR framework of \textcite{chernozhukov_hansen_2005} uses excluded instruments to connect endogenous treatment variation to structural outcome quantiles.

Structural IVQR represents potential outcomes through a quantile function and a latent structural rank. For a treatment state $d$, the structural rank $U_d\in(0,1)$ is the latent quantile index that locates an individual within the conditional distribution of the potential outcome $Y_d$ given observed controls. Because the structural quantile function is increasing in this rank, a larger value of $U_d$ corresponds to a higher position in the potential-outcome distribution. The rank may be interpreted as summarizing unobserved heterogeneity that determines where an observationally similar individual lies in that distribution. Instrument exogeneity restricts the relationship between the instruments and these structural ranks conditional on controls, while treatment may still depend on latent variables related to the outcome. Rank invariance assumes that an individual's rank is unchanged across treatment states; the weaker rank-similarity condition allows ranks to change but requires their conditional distributions to agree across those states. Under the maintained selection, exogeneity, and rank conditions, the probability that the observed outcome lies below the true structural quantile, conditional on controls and instruments, equals the quantile index \parencite{chernozhukov_hansen_2005,chernozhukov_hansen_wuthrich_2017}. This conditional restriction links the potential-outcome model to estimable instrument moments. Validity permits the instrument to enter those moments; relevance determines whether instrument-induced treatment variation separates competing structural coefficients. Moment validity alone does not establish point identification. A unique solution additionally requires the population mapping from candidate structural quantiles to the restrictions to distinguish the truth from alternatives. That mapping is quantile specific and density weighted, so a mean first-stage statistic is informative about relevance but does not completely characterize IVQR identification at a particular quantile. This distinction is essential empirically: an instrument can satisfy the maintained exogeneity condition yet generate too little treatment variation to locate a structural quantile effect precisely.

When identifying information is strong, the population moment moves sufficiently as the candidate treatment coefficient moves away from its true value, supporting regular point-estimator approximations. With weak information, the true value may remain unique while nearby candidates generate almost indistinguishable moments. Partial identification instead permits more than one population solution, and failed identification leaves some candidates or directions indistinguishable. \textcite{chernozhukov_hansen_2008} develop inference that remains asymptotically valid without imposing the strong-identification conditions needed for conventional estimate-centered inference. Their approach fixes a candidate value, estimates or profiles the nuisance parameters under that candidate, and tests whether the resulting instrument-based quantile moments are compatible with zero. Inverting these candidate-specific tests retains every non-rejected value and can therefore produce asymmetric, disconnected, or very wide confidence regions. Weak-identification robustness concerns the null reference distribution; it does not guarantee that false candidates will be rejected. Informativeness and finite-sample calibration remain distinct: a robust confidence region can honestly reveal little about the parameter, and an asymptotic test can still reject too often or too rarely in a finite sample. The exact conditional procedure of \textcite{chernozhukov_hansen_jansson_2009} uses a different simulated reference distribution and is not the method implemented here.

The IVQR inference used in this thesis also belongs to a broader literature on inference under weak instrumental-variable identification. \textcite{anderson_rubin_1949} provide the classical foundation for testing structural parameter values through restrictions that remain meaningful without relying on a precisely estimated first stage. \textcite{staiger_stock_1997} formalize weak-instrument asymptotics in linear instrumental-variables models and show why conventional estimate-centered approximations can become unreliable when instruments are only weakly correlated with endogenous regressors. \textcite{stock_wright_2000} extend weak-identification analysis to generalized method-of-moments settings and develop confidence procedures that remain valid when parameters are weakly identified. Alternative robust procedures include the pivotal testing approach of \textcite{kleibergen_2002} and the conditional likelihood-ratio method of \textcite{moreira_2003}. \textcite{andrews_stock_sun_2019} provide a modern review of this literature and its implications for applied work. These results motivate the general distinction between conventional estimate-centered inference and candidate-specific procedures that remain valid under weak identification. They do not, however, provide a complete characterization of identification in IVQR. In particular, linear-IV first-stage diagnostics and threshold rules do not directly measure quantile-specific IVQR identification, whose population moments are nonlinear and depend on the distribution of the outcome around the quantile of interest. The thesis therefore uses first-stage statistics only to describe the designed ordering of excluded-instrument relevance, while the IVQR-specific weak-identification-robust inference follows \textcite{chernozhukov_hansen_2008}.

High-dimensional controls create a second challenge because estimating the outcome and instrument nuisances may require many coefficients relative to the sample size. Regularization exploits approximate sparsity: a comparatively small set of variables captures most predictive content, while the remaining coefficients are zero or sufficiently small. Shrinkage and selection can nevertheless contaminate a target moment. Double machine learning addresses this problem with a Neyman-orthogonal score whose population expectation has zero first derivative with respect to local nuisance perturbations at the truth \parencite{chernozhukov_et_al_2018}. Orthogonality therefore removes the leading sensitivity to sufficiently small nuisance-estimation errors, not the errors themselves; the learners must still satisfy suitable convergence and regularity conditions. Sample splitting estimates nuisances on one subsample and evaluates the score on another, while cross-fitting rotates these roles across folds. These concepts should not be conflated: orthogonality is a property of the score, whereas cross-fitting organizes data use and limits overfitting dependence. Nor does DML strengthen an instrument. It addresses high-dimensional nuisance estimation, while candidate-specific robust inference addresses irregular inference under limited identification.

\textcite{chen_huang_tien_2021} join these two strands. For each candidate, they profile a high-dimensional quantile regression, construct a density-weighted instrument residual, form a Neyman-orthogonal IVQR score, estimate the nuisances with regularization, and evaluate candidates by GMM and weak-identification-robust test inversion. Although they discuss cross-fitting as part of the broader DML framework, the executable DML-IVQR algorithm preserved in this thesis estimates nuisances and evaluates the score on the same sample. Their principal Monte Carlo design holds instrument relevance fixed while comparing Oracle-GMM, Full-GMM, and DML-IVQR across several quantiles. The remaining empirical question is how point-estimation error, confidence-set informativeness, power, and finite-sample calibration respond when excluded-instrument relevance is systematically weakened within that otherwise preserved design. This thesis supplies that controlled stress test; it does not propose a new structural model, score, estimator, identification theorem, or finite-sample correction. \textcite{jin_sun_2025} study a related nonparametric high-dimensional IVQR problem for unconditional quantile treatment effects with automatic DML and cross-fitting, but their target and design differ from the controlled relevance experiment considered here. Chapter~3 states the exact structural model, nuisance construction, estimators, and inference procedure used in the simulations.
